% Part: incompleteness
% Chapter: representability-in-q
% Section: composition-representable

\documentclass[../../../include/open-logic-section]{subfiles}

\begin{document}

\olfileid{inc}{req}{cmp}
\olsection{సంయుక్తం~$\Th{Q}$లో ప్రాతినిధ్యం చేయదగినది}

$h$ను ఇలా నిర్వచించారని అనుకుందాం:
\[
h(x_0,\dots,x_{l-1}) = f(g_0(x_0,\dots,x_{l-1}), \dots,
g_{k-1}(x_0,\dots,x_{l-1})).
\]
ఇక్కడ $f$, $g_0$,~\dots,~$g_{k-1}$ ప్రమేయాలకు వరుసగా ప్రాతినిధ్యం చేసే
\tetoken{సూత్రాలు}{!!{formula}s} $!A_f,!A_{g_0},\dots,!A_{g_{k-1}}$లను ఇప్పటికే కనుగొన్నాం.
$h$కు ప్రాతినిధ్యం చేసే \tetoken{ఒక సూత్రం}{!!a{formula}}~$!A_h$ను కనుగొనాలి.

అన్ని ప్రమేయాలూ $1$-స్థానికమైన $h(x)=f(g(x))$ అనే సరళ సందర్భంతో
మొదలుపెడదాం. $!A_f(y,z)$ అనేది~$f$కు, $!A_g(x,y)$ అనేది~$g$కు
ప్రాతినిధ్యం చేస్తే, $h$కు ప్రాతినిధ్యం చేసే \tetoken{ఒక సూత్రం}{!!a{formula}}~$!A_h(x,z)$
కావాలి. $z=f(y)$, $y=g(x)$ రెండూ నెరవేర్చే ఒక~$y$ ఉన్నప్పుడు, అప్పుడు మాత్రమే
$h(x)=z$ అని గమనించండి. ($h(x)=z$ అయితే $g(x)$ అలాంటి~$y$; అలాంటి~$y$
ఉంటే $y=g(x)$, $z=f(y)$ కనుక $z=f(g(x))$.) అందువల్ల
$\lexists[y][(!A_g(x,y)\land !A_f(y,z))]$ను $!A_h(x,z)$కు మంచి అభ్యర్థిగా
తీసుకోవచ్చు. $\Th{Q}$ సంబంధిత \tetoken{సూత్రాలను}{!!{formula}s} నిరూపిస్తుందని మాత్రమే
ధృవీకరించాలి.

\begin{prop}
\ollabel{prop:rep1}
$h(n)=m$ అయితే $\Th{Q}\Proves !A_h(\num{n},\num{m})$.
\end{prop}

\begin{proof}
$h(n)=m$, అంటే $f(g(n))=m$ అనుకుందాం. $k=g(n)$గా తీసుకోండి. అప్పుడు
\begin{align*}
  \Th{Q} & \Proves !A_g(\num{n}, \num{k})
  \intertext{$!A_g$ అనేది~$g$కు ప్రాతినిధ్యం చేస్తుంది కనుక; అలాగే}
  \Th{Q} & \Proves !A_f(\num{k}, \num{m})
  \intertext{$!A_f$ అనేది~$f$కు ప్రాతినిధ్యం చేస్తుంది కనుక. అందువల్ల}
  \Th{Q} & \Proves !A_g(\num{n}, \num{k}) \land !A_f(\num{k}, \num{m})
  \intertext{మరియు పర్యవసానంగా}
  \Th{Q} & \Proves \lexists[y][(!A_g(\num{n}, y) \land !A_f(y, \num{m}))],
\end{align*}
అంటే $\Th{Q}\Proves !A_h(\num{n},\num{m})$.
\end{proof}

\begin{prop}
\ollabel{prop:rep2}
$h(n)=m$ అయితే $\Th{Q}\Proves\lforall[z][(!A_h(\num{n},z)\lif
z=\num{m})]$.
\end{prop}

\begin{proof}
$h(n)=m$, అంటే $f(g(n))=m$ అనుకుందాం. $k=g(n)$గా తీసుకోండి. అప్పుడు
\begin{align*}
  \Th{Q} & \Proves \lforall[y][(!A_g(\num{n}, y) \lif \eq[y][\num{k}])]
  \intertext{$!A_g$ అనేది~$g$కు ప్రాతినిధ్యం చేస్తుంది కనుక; అలాగే}
  \Th{Q} & \Proves \lforall[z][(!A_f(\num{k}, z) \lif \eq[z][\num{m}])]
  \intertext{$!A_f$ అనేది~$f$కు ప్రాతినిధ్యం చేస్తుంది కనుక. కొద్దిపాటి
    తర్కాన్ని వాడి కిందిదీ వస్తుందని చూపవచ్చు:}
  \Th{Q} & \Proves \lforall[z][(\lexists[y][(!A_g(\num{n}, y) \land
      !A_f(y, z))] \lif \eq[z][\num{m}])].
\end{align*}
అంటే $\Th{Q}\Proves\lforall[y][(!A_h(\num n,y)\lif\eq[y][\num m])]$.
\end{proof}

$f$, $g_i$ల స్థానసంఖ్య $1$కంటే ఎక్కువగా ఉండే సంక్లిష్ట సందర్భంలోనూ ఇదే
ఆలోచన పనిచేస్తుంది.

\begin{prop}
\ollabel{prop:rep-composition}
$!A_f(y_0,\dots,y_{k-1},z)$ అనేది~$f(y_0,\dots,y_{k-1})$కు~$\Th{Q}$లో
ప్రాతినిధ్యం చేస్తుందనీ, $!A_{g_i}(x_0,\dots,x_{l-1},y)$ అనేది
$g_i(x_0,\dots,x_{l-1})$కు~$\Th{Q}$లో ప్రాతినిధ్యం చేస్తుందనీ అనుకుందాం.
అప్పుడు
\begin{multline*}
  \lexists[y_0][\dots\lexists[y_{k-1}][(!A_{g_0}(x_0,\dots,x_{l-1},y_0) \land
      \dots \land {}\\
  !A_{g_{k-1}}(x_0,\dots,x_{l-1},y_{k-1}) \land
  !A_f(y_0,\dots,y_{k-1},z))]\dots]
\end{multline*}
అనే సూత్రం
\[
h(x_0, \dots, x_{l-1}) = f(g_0(x_0, \dots, x_{l-1}), \dots, g_{k-1}(x_0,
\dots, x_{l-1})).
\]
కు ప్రాతినిధ్యం చేస్తుంది.
\sourcecorrection{OLTEREQ-005}{ఆంగ్ల మూలంలోని సాధారణ సంయుక్త ప్రాతినిధ్య సూత్రంలో పరిమాణక గూళ్ల మూసివేతకు ముందే సంధానంలోని చివరి రెండు సూత్రాలు బయటపడి, కుండలీకరణ అసమతుల్యమైంది. ప్రతిపాదనకు ముందు ఇచ్చిన ఏకస్థానిక నమూనా నిర్దేశించేలా అన్ని $g_i$ సూత్రాలనూ $A_f$నూ ఒకే సంధానంలో, $y_0,\dots,y_{k-1}$ అస్తిత్వ పరిమాణకాల పరిధిలో ఉంచాం.}
\end{prop}

\begin{proof}
సాధన.
\end{proof}

\begin{prob}
\olref[inc][req][cmp]{prop:rep1},
\olref[inc][req][cmp]{prop:rep2} నిరూపణలను మార్గదర్శకంగా తీసుకుని,
\olref[inc][req][cmp]{prop:rep-composition} నిరూపణను వివరంగా చేయండి.
\sourcecorrection{OLTEREQ-006}{ఆంగ్ల మూలంలోని సమస్య మార్గదర్శకాలుగా $\olref[inc][req][cmp]{prop:rep2}$నే రెండుసార్లు సూచించింది. వెంటనే ముందున్న రెండు వేర్వేరు నిరూపణలను ఉద్దేశించినందున మొదటి సూచనను $\olref[inc][req][cmp]{prop:rep1}$గా పునరుద్ధరించాం.}
\end{prob}

\end{document}
